% Part: computability
% Chapter: computability-theory
% Section: no-universal-function

\documentclass[../../../include/open-logic-section]{subfiles}

\begin{document}

\olfileid[pt]{cmp}{thy}{nou}
\olsection{Inexistência de uma função computável universal}

Embora exista uma função parcialmente computável que é total para as
funções parcialmente computáveis, não existe uma função computável
total que seja universal para as funções computáveis totais.

\begin{thm}
\ollabel{thm:no-univ}
Não existe uma função computável universal. Em outras palavras,
qualquer função $\fn{Un}'(k, x)$ tal que, se $f(x)$ é uma função
computável total, então existe um número natural~$k$ tal que
$f(x) = \fn{Un}'(k,x)$ para todo~$x$, não é computável.
\end{thm}

\begin{proof}
A demonstração é uma diagonalização simples: se $\fn{Un}'(k,x)$ fosse
total e computável, então
\[
d(x) = \fn{Un}'(x, x) + 1
\]
também seria total e computável. Contudo, por definição, $d(k)$ não é
igual a $\fn{Un}'(k,k)$. Portanto, para todo $k$, os valores de
$d(x)$ e~$\fn{Un}'(k, x)$ diferem para pelo menos um~$x$, a saber,
$x = k$.
\end{proof}

\begin{explain}
A \olref[uni]{thm:univ-comp} acima mostra que podemos contornar esse
argumento de diagonalização, mas apenas ao custo de permitir que a
função universal seja parcial. Isto é, $\fn{Un}$ é universal para as
funções computáveis totais; ela apenas não é total. O argumento de
diagonalização não funciona no caso parcial.
\end{explain}

\begin{prob}
  Para entender por que o argumento de diagonalização na demonstração
  de \olref{thm:no-univ} não funciona no caso parcial, considere a
  função $f(x) \simeq \fn{Un}(x,x)+1$. Ela é parcialmente computável?
  Se for, tem um índice~$e$, isto é,
  $f(x) \simeq \fn{Un}(e,x)$. O que você pode dizer sobre~$f(e)$?
\end{prob}


\end{document}
